\FifthArticle{10}

It remains only to show that, if the point $P$ representing a root $=p$\SourceLeaf{090} lies on two or more lines $T=0$ at once, the square of $x-p$, or the higher power corresponding to the number of these concurrent lines, will be contained in $X$ as a factor. The proof rests on the following propositions.

Introduce another indeterminate quantity $z$ in place of $x$ by putting \mbox{$x=z+p$}. Let this substitution turn $X$ into $Z$, so that $Z$ is a function of $z$ of the same degree as $X$ in $x$, but without a constant term. Arrange it in ascending powers of $z$, and let its lowest nonvanishing term be
\[
=Kz^m\quad\mbox{and}\quad Z=Kz^m(1+\zeta)
\]
where $\zeta$ will have the form $Lz+L'zz+L''z^3+\mbox{etc.}+\frac{1}{K}z^{n-m}$. Finally put
\[
z=s(\cos\psi+i\sin\psi)
\]

The real and imaginary parts of $z$ express the position of any indeterminate point of the plane as rectangular coordinates, and the quantities $s$, $\psi$ as polar coordinates, relative to the point $P$, in exactly the same way as the parts of $x$ and the quantities $r$, $\varphi$ designate the relative position with respect to the original origin. The combination of a fixed value of $s$ with all values of $\psi$ over a range of $360^0$ thus represents the points of a circumference with centre $P$ and radius $=s$.

If one now puts $K=k(\cos\kappa+i\sin\kappa)$, and consequently
\[
Kz^m=ks^m(\cos(m\psi+\kappa)+i\sin(m\psi+\kappa))
\]
then, for an infinitely small $s$, the quantity $\zeta$, which is of at least the same order as $s$, may be neglected beside $1$, and it is therefore permissible to put
\[
T=ks^m\cos(m\psi+\kappa)
\]
It follows that, as $\psi$ increases by $360^0$, the sign of $T$ is positive on $m$ pieces of the circumference and negative on just as many alternating with them; or that $T$ takes the value $=0$ at $2m$ points, namely for $\psi=\frac{1}{m}(\kappa-90^0)$, $\frac{1}{m}(\kappa+90^0)$, $\frac{1}{m}(\kappa+270^0)$, etc. Thus a total of $2m$ lines on which $T=0$ issue from $P$\SourceLeaf{091}. Or, if they are joined in pairs so that each line at which, as $\psi$ increases, the sign passes from $-$ to $+$ is regarded, together with the next one at which the opposite transition occurs, as the boundary of a region with positive $T$, then in all $m$ such boundary lines meet at $P$.

On the other hand, it is clear that, just as $Z$ is identically divisible by $z^m$ and by no higher power of $z$, $X$ will contain the factor $(x-p)^m$, but no higher power of $x-p$. Thus, whenever $p$ denotes any root of the equation $X=0$, the exponent of the highest power of $x-p$ by which $X$ is divisible is equal to the number of boundary lines of regions with positive $T$ meeting \emph{at} $P$, or, what is the same, to the number of such regions meeting \emph{at} $P$.

Moreover, it is easy to give the proof a form entirely independent of any introduction of infinitely small quantities, and quite analogous to the chain of reasoning in Articles 3 and 4. Namely, one can establish a value of $s$ for which, as for every smaller value, the entire cycle of all values of $\psi$ presents the same alternating succession of $m$ pieces with positive $T$ and just as many with a negative value. The positive root of the equation
\[
0=m\surd\frac{1}{2}-(m+1)ls-(m+2)l'ss-(m+3)l''s^3-\mbox{etc.}
\]
has this property. Here $l$, $l'$, $l''$, etc. denote the positive square roots of the norms of the complex quantities $L$, $L'$, $L''$, etc.; alternatively, put
\[
\begin{array}{rl}
L&=l\,(\cos\lambda+i\sin\lambda)\\
L'&=l'(\cos\lambda'+i\sin\lambda')\\
L''&=l''(\cos\lambda''+i\sin\lambda'')\ \mbox{etc.}
\end{array}
\]
I believe, however, that I may pass over the very easy development of this proposition here.

Finally, it may be observed that the proof in the paper of 1799 required the consideration of \emph{two} systems of lines, one containing the lines where $T=0$, the other those where $U=0$, whereas in our present procedure the consideration of one system\SourceLeaf{092} has sufficed. For this I chose the system of boundary lines of the regions with positive $T$, but consideration of the boundary lines of the regions with positive (or negative) $U$ could have served the same purpose equally well.

